% TOP_PROBLEM: {"id": 3394, "title": "Sums of independent random variables with unbounded variance", "category_id": 15, "category": {"id": 15, "name": "computer_science", "display_name": "Computer Science", "description": "Computational complexity, algorithms, and theoretical CS.", "slug": "computer-science", "order_index": 15, "created_at": "2026-07-31T15:26:25.671Z"}, "status": "open", "problem_number": "OPG-36887", "set_id": 12}
% FIRST_POSTED: 2026-10-01
% ENTRY_KIND: statement_only
% UnsolvedMath ID 3394; source catalogue code OPG-36887
% !TeX program = lualatex
% Definition and sources only; this file is not an LLM solution attempt.
\documentclass[11pt,a4paper]{article}
\usepackage[margin=25mm]{geometry}
\usepackage{fontspec}
\setmainfont{lmroman10-regular.otf}[BoldFont=lmroman10-bold.otf,
 ItalicFont=lmroman10-italic.otf,BoldItalicFont=lmroman10-bolditalic.otf]
\usepackage{amsmath,amssymb,mathtools}
\usepackage{unicode-math}
\setmathfont{latinmodern-math.otf}
\AtBeginDocument{\let\setminus\smallsetminus}
\usepackage{xurl,hyperref}
\hypersetup{colorlinks=true,urlcolor=blue}
\setcounter{secnumdepth}{0}
\setlength{\parindent}{0pt}
\setlength{\parskip}{5pt}
\setlength{\emergencystretch}{3em}
\begin{document}
\section*{Sums of independent random variables with unbounded variance}\label{3394}
{\small UnsolvedMath ID 3394 · OPG-36887 · Computer Science · OpenGarden}
\subsection{Definitions and mathematical statement}
Let $n\ge1$, $\mu>0$ and $\delta>0$. Let $X_1,\ldots,X_n$ be
independent nonnegative real random variables on a probability space,
with finite expectations satisfying $\mathbb E X_i\le\mu$ for each
$i$. They need not have identical distributions, and their variances
may be infinite. Write $S=\sum_{i=1}^n X_i$, so
$\mathbb E S=\sum_i\mathbb E X_i$ is finite.

The conjectured uniform lower bound is
\[
 \mathbb P(S-\mathbb E S<\delta\mu)
            \ge \min\left\{\frac{\delta}{1+\delta},\ e^{-1}\right\}.
\]
The event uses a strict inequality, which matters when the variables
have atoms. The centering is the actual expectation of the sum,
rather than the upper bound $n\mu$. The proposed right-hand side
depends only on $\delta$, independently of $n$, $\mu$, and all
the distributions. Dividing each $X_i$ by $\mu$ gives the equivalent
normalization $\mu=1$; requiring $\mu>0$ avoids a degenerate strict
event at zero.
Establish the inequality for every admissible family and every
$\delta>0$, or give a fully specified independent family and
parameter violating it. An assumption of bounded variance or
bounded support would restrict the stated problem.
\subsection{Short English statement}
How small can the probability be that a sum of independent nonnegative variables lies less than a prescribed amount above its own mean?
\subsection{Sources}
\begin{itemize}
\item[S1] ULAM.AI and the UnsolvedMath Contributors, supplied problems.json, record 3394 (OPG-36887), OpenGarden collection. Catalogue identity and source transcription; CC BY 4.0.
\url{https://huggingface.co/datasets/ulamai/UnsolvedMath}.
\item[S2] Open Problem Garden, Sums of independent random variables with unbounded variance. Original problem and discussion; the mathematical formulation below is expanded from the supplied source record.
\url{http://www.openproblemgarden.org/op/sums_of_independent_random_variables_with_unbounded_variance}.
\end{itemize}
\end{document}
